% Chapter 2: Literature Review - State of the Art in Ontology-Grounded RAG Systems
% Author: Charles Watson Ndethi Kibaki
% Date: August 2025

\section{Introduction to the State-of-the-Art Analysis}

This literature review presents a comprehensive State-of-the-Art (SotA) analysis of Ontology-Grounded Retrieval-Augmented Generation (OG-RAG) systems, structured around a three-part argument that traces the evolution, current capabilities, and future potential of this rapidly advancing field. Following established SotA methodology, this review addresses: (1) \textit{"This is where we are now"} - the contemporary paradigm of OG-RAG systems as of 2024-2025; (2) \textit{"This is how we got here"} - the historical evolution and pivotal breakthroughs that shaped current understanding; and (3) \textit{"This is where we could go next"} - emerging research directions and critical gaps that define future opportunities.

The temporal scope of this review is deliberately anchored around December 2024, which represents a watershed moment in the field with Microsoft Research's breakthrough demonstration of 55\% improvement in factual recall and 40\% enhancement in response correctness through ontology-grounded approaches \cite{chen2024og}. This landmark achievement fundamentally shifted the research paradigm from conventional vector-similarity RAG toward structured knowledge integration, establishing the foundation for contemporary OG-RAG systems.

\section{Part I: Current State of Ontology-Grounded RAG (2024-2025)}
\subsection{Defining the Contemporary Paradigm}

The current state of Ontology-Grounded Retrieval-Augmented Generation represents a paradigm shift from the limitations of conventional RAG systems that rely primarily on vector similarity for context retrieval. The pivotal moment came in December 2024 with Microsoft Research's publication of "Ontology-Grounded Retrieval-Augmented Generation for Large Language Models" \cite{chen2024og}, which demonstrated unprecedented improvements in factual accuracy and response correctness through explicit ontological grounding.

This breakthrough established the contemporary understanding that effective RAG systems require more than semantic similarity; they demand structured knowledge representations that preserve complex relationships, hierarchical dependencies, and domain-specific constraints. The Microsoft OG-RAG framework introduced hypergraph representations where each hyperedge encapsulates clusters of factual knowledge precisely grounded in domain-specific ontologies \cite{chen2024og}. This architectural innovation enables retrieval of minimal sets of hyperedges that collectively form conceptually grounded contexts, moving beyond the fragmented text chunks characteristic of traditional RAG approaches.

The quantitative evidence supporting this paradigm shift is compelling. Across various Large Language Models, OG-RAG systems consistently demonstrate 55\% improvements in recall of accurate facts, 40\% enhancements in overall response correctness, 30\% faster attribution of responses to context, and 27\% increases in fact-based reasoning accuracy \cite{chen2024og}. These performance gains represent not incremental improvements but fundamental advances in how knowledge can be structured and leveraged for generation tasks.

\subsection{Current Technical Landscape and Dominant Frameworks}

The contemporary technical landscape of OG-RAG systems is characterized by several dominant architectural paradigms, each addressing different aspects of the knowledge integration challenge. These frameworks represent the state-of-the-art in 2024-2025 and define the current capabilities and limitations of the field.

\subsubsection{Hypergraph-Based Architectures}

Hypergraph-based systems represent the most sophisticated current approach to ontology-grounded retrieval. The Microsoft OG-RAG framework leads this category, employing hypergraphs where individual hyperedges represent clusters of factual knowledge anchored in domain ontologies \cite{chen2024og}. The HyperGraphRAG extensions \cite{wang2024hypergraphrag} further advance this approach by enabling modeling of n-ary relationships that cannot be captured in traditional binary graph structures.

These hypergraph architectures address the fundamental limitation of conventional RAG systems that struggle with complex relational knowledge. By representing knowledge as hyperedges rather than simple edges, these systems can capture multi-entity relationships, conditional dependencies, and hierarchical structures that are essential for domain-specific reasoning. The optimization algorithms employed for hyperedge selection ensure that retrieved contexts maintain logical coherence while minimizing redundancy \cite{chen2024og}.

\subsubsection{Graph Neural Network Integration}

GNN-RAG frameworks represent another major architectural paradigm, exemplified by the work of Mavromatis and Karypis \cite{mavromatis2024gnnrag}, which introduces dual-stage reasoning mechanisms. These systems achieve 8.9-15.5\% F1 score improvements on multi-hop question answering tasks by explicitly modeling graph structures through Graph Neural Networks integrated with transformer architectures.

The GNN-RAG approach addresses the challenge of reasoning over complex graph structures by employing specialized graph encoders that capture topological features alongside textual content. This dual encoding strategy enables the system to leverage both the semantic content of nodes and the structural relationships encoded in the graph topology. The integration with transformer models allows for end-to-end training while preserving the inductive biases inherent in graph structures \cite{mavromatis2024gnnrag}.

G-Retriever \cite{he2024gretriever}, presented at NeurIPS 2024, represents the first RAG system designed specifically for general textual graphs, demonstrating the generalizability of GNN-based approaches beyond domain-specific knowledge graphs. This framework enables flexible graph representation learning that adapts to diverse graph structures and content types.

\subsubsection{Multimodal and Neurobiologically-Inspired Approaches}

The integration of multimodal capabilities represents an emerging frontier in OG-RAG systems. GraphVis \cite{zhang2024graphvis}, also presented at NeurIPS 2024, demonstrates 11.1\% average improvements across multiple benchmarks by incorporating visual and textual graph information through curriculum fine-tuning strategies. This multimodal integration is particularly relevant for cultural heritage applications where visual artifacts and textual descriptions are interconnected.

HippoRAG \cite{bai2024hipporag} introduces a neurobiologically-inspired framework that mimics human memory mechanisms through Personalized PageRank algorithms. This approach demonstrates how insights from cognitive science can inform the design of retrieval mechanisms, particularly for scenarios requiring episodic memory-like capabilities. The system's ability to model memory consolidation and retrieval processes offers unique advantages for applications requiring long-term knowledge retention and contextual recall.

\subsubsection{Production Systems and Industry Adoption}

The transition from research prototypes to production systems marks a critical indicator of technological maturity. Microsoft GraphRAG has achieved enterprise-level implementation with demonstrated scalability across large-scale knowledge graphs \cite{edge2024graphrag}. The system's integration with Azure Cognitive Services provides a pathway for real-world deployment while maintaining the sophisticated reasoning capabilities demonstrated in research settings.

Neo4j's comprehensive GraphRAG support represents another milestone in industry adoption \cite{neo4j2024graphrag}. The platform's integration of Cypher query language with vector indices enables hybrid retrieval strategies that combine graph traversal with semantic similarity. This hybrid approach addresses practical deployment challenges while preserving the structured reasoning capabilities that define OG-RAG systems.

Cost optimization has emerged as a critical practical consideration, addressed by innovations such as LazyGraphRAG \cite{guo2024lazygraphrag} and the 3B Phi-3 implementation in Triplex \cite{zhang2024triplex}. These developments demonstrate that sophisticated ontology-grounded reasoning can be achieved with smaller, more efficient models, making the technology accessible for resource-constrained applications.

The open-source ecosystem has also matured significantly, with frameworks like LightRAG \cite{guo2024lightrag} and RAPTOR \cite{sarthi2024raptor} providing accessible implementations of graph-based retrieval mechanisms. These community-driven initiatives ensure that OG-RAG capabilities are not limited to large technology companies but can be adapted and extended by researchers and practitioners worldwide.

\section{Part II: Historical Evolution and Pivotal Breakthroughs}

\subsection{Pre-OG-RAG Foundations and Early Limitations}

The historical trajectory toward ontology-grounded RAG systems begins with the recognition of fundamental limitations in traditional RAG approaches that emerged in the early 2020s. The original RAG framework \cite{lewis2020retrieval} demonstrated the potential for augmenting language models with external knowledge but revealed critical shortcomings when applied to complex, structured domains.

Traditional RAG systems operated under the assumption that vector similarity in embedding spaces would adequately capture semantic relationships. However, empirical studies revealed that this approach failed to preserve logical dependencies, hierarchical structures, and domain-specific constraints \cite{petroni2021kilt}. The reliance on dense vector representations obscured the explicit relationships between entities, leading to context fragmentation and logical inconsistencies in generated responses.

Early attempts to address these limitations focused on improving embedding quality and retrieval strategies without fundamentally questioning the vector-similarity paradigm. BM25 hybrid approaches \cite{karpukhin2020dense} and dense passage retrieval optimizations \cite{qu2021rocketqa} achieved incremental improvements but could not overcome the inherent limitation of treating knowledge as isolated text chunks rather than interconnected conceptual structures.

The recognition that knowledge graphs could provide the structured foundation missing from traditional RAG emerged gradually through multiple research streams. Work on knowledge graph completion \cite{sun2019rotate} and graph neural networks for reasoning \cite{hamilton2017inductive} established the technical foundations for graph-based knowledge representation. Simultaneously, research on knowledge-grounded dialogue systems \cite{moon2019opendialkg} demonstrated the importance of structured knowledge for maintaining contextual coherence in generation tasks.

\subsection{Breakthrough Period: The Path to Ontological Grounding}

The transition from traditional RAG to ontology-grounded approaches can be traced through several pivotal breakthroughs that fundamentally changed the field's understanding of knowledge integration.

\subsubsection{Early Graph Integration Attempts}

The first significant step toward ontology-grounded RAG came with KG-RAG implementations that attempted to integrate knowledge graphs directly into retrieval pipelines \cite{yu2022kg}. These early systems demonstrated the potential for structured knowledge to improve factual accuracy but struggled with scalability and integration challenges. The approaches were often hybrid, maintaining separate vector and graph retrieval systems without true integration.

ERNIE \cite{zhang2019ernie} and KnowBERT \cite{peters2019knowledge} represented important precursors by integrating entity knowledge directly into language model pre-training. While not RAG systems per se, these approaches demonstrated that explicit knowledge integration could significantly improve model performance on knowledge-intensive tasks. The success of these knowledge-enhanced language models provided empirical evidence for the value of structured knowledge integration.

\subsubsection{Graph Neural Network Foundations}

The development of Graph Neural Networks specifically for textual graphs marked a crucial turning point. GraphSAINT \cite{zeng2019graphsaint} and subsequent architectures demonstrated that GNNs could effectively encode both textual content and graph structure simultaneously. This capability addressed the fundamental challenge of preserving structural relationships while maintaining compatibility with transformer-based language models.

The introduction of heterogeneous graph neural networks \cite{wang2019heterogeneous} enabled modeling of complex knowledge graphs with multiple entity and relation types. This advancement was essential for real-world applications where knowledge graphs contain diverse entity categories and relationship hierarchies that cannot be captured by homogeneous graph models.

\subsubsection{The Ontology Integration Breakthrough}

The conceptual breakthrough that led to contemporary OG-RAG systems came with the recognition that ontologies could provide the formal foundation for grounding retrieval processes. Unlike ad hoc knowledge graphs, ontologies offer formal semantics, logical consistency, and domain-specific vocabularies that can guide both knowledge representation and retrieval strategies.

The work on ontology-grounded dialogue systems \cite{he2022ontology} demonstrated that formal ontological structures could significantly improve the coherence and accuracy of knowledge-grounded generation. This research established the theoretical foundation for using ontologies not merely as knowledge sources but as structural frameworks for organizing and accessing knowledge.

The development of hybrid reasoning systems that combined symbolic ontological reasoning with neural generation \cite{bosselut2019comet} represented another crucial advancement. These systems demonstrated that symbolic and neural approaches could be effectively integrated, with ontologies providing logical structure and neural models providing flexible generation capabilities.

\subsubsection{The Microsoft OG-RAG Watershed Moment}

The December 2024 publication of Microsoft's OG-RAG framework \cite{chen2024og} represents the culmination of these historical developments and establishes the contemporary paradigm. This work synthesized insights from knowledge graph reasoning, graph neural networks, and ontological engineering into a unified framework that achieved unprecedented performance improvements.

The key innovation was the introduction of hypergraph representations where hyperedges are explicitly grounded in domain ontologies. This approach resolved the long-standing challenge of maintaining logical coherence while enabling flexible retrieval strategies. The optimization algorithms for hyperedge selection ensured that retrieved contexts preserved both semantic relevance and logical consistency.

The quantitative results—55\% improvement in factual recall and 40\% enhancement in response correctness—provided compelling empirical evidence for the superiority of ontology-grounded approaches. These results were not merely incremental improvements but represented a qualitative shift in system capabilities, demonstrating that explicit ontological grounding could overcome fundamental limitations of vector-similarity approaches.

\subsection{Convergent Research Streams}

The path to contemporary OG-RAG systems involved the convergence of several previously independent research streams, each contributing essential capabilities to the final synthesis.

Research on automated ontology construction \cite{fernandez2019ontology} provided the tools necessary for building domain-specific ontologies at scale. Without these automated approaches, the manual construction of comprehensive ontologies would have been prohibitively expensive and time-consuming.

Advances in graph database technologies \cite{angles2018foundations} enabled the storage and querying of large-scale knowledge graphs with the performance characteristics required for real-time retrieval. The development of specialized query languages and indexing strategies made it feasible to deploy graph-based retrieval systems in production environments.

The emergence of transformer architectures with explicit support for structured inputs \cite{shaw2018self} provided the neural foundation for integrating graph-structured knowledge with language models. The ability to encode positional and structural information within transformer attention mechanisms was essential for maintaining the benefits of structured knowledge while leveraging the generative capabilities of large language models.

\section{Part III: Future Research Directions and Critical Gaps}

\subsection{Critical Research Gaps and Scalability Challenges}

Despite the significant advances in OG-RAG systems, several critical research gaps limit the widespread adoption and effectiveness of these approaches, particularly for specialized domains and low-resource contexts.

\subsubsection{Resource-Intensive Ontology Construction}

The most significant barrier to OG-RAG adoption remains the resource-intensive nature of ontology construction. Current approaches require substantial domain expertise and manual effort to create comprehensive, accurate ontologies \cite{fernandez2019ontology}. For specialized domains like cultural heritage or low-resource languages, the scarcity of domain experts compounds this challenge. Automated ontology learning techniques show promise but struggle with the nuanced cultural and contextual knowledge required for faithful representation of complex domains.

Recent work on LLM-based ontology population \cite{agarwal2024llm} demonstrates 61.05\% F1 score improvements over traditional span-based methods, suggesting that large language models could assist in ontology construction. However, these approaches still require significant human oversight and validation, particularly for culturally sensitive domains where accuracy and authenticity are paramount.

\subsubsection{Cross-Domain Generalization and Transfer Learning}

Current OG-RAG systems typically require domain-specific ontologies and cannot easily transfer knowledge across domains. This limitation is particularly problematic for applications that span multiple domains or require integration of diverse knowledge sources. The development of foundational ontologies that can serve as bridges between domain-specific knowledge representations remains an open challenge.

The work on universal graph neural networks \cite{you2021graph} provides some foundation for cross-domain transfer, but effective transfer of ontological knowledge requires preserving semantic relationships across domain boundaries. Current approaches struggle with the alignment of concepts and relationships between ontologies, limiting their ability to leverage knowledge learned in one domain for applications in related domains.

\subsubsection{Cultural Knowledge Representation and Ethical Frameworks}

The application of OG-RAG systems to cultural heritage and indigenous knowledge raises significant ethical and methodological challenges that current research has not adequately addressed. The representation of cultural knowledge in formal ontologies risks imposing Western epistemological frameworks on knowledge systems that may not align with such structures \cite{christie2019indigenous}.

The principles of Indigenous Data Sovereignty, including OCAP (Ownership, Control, Access, Possession) and FPIC (Free, Prior, and Informed Consent), require fundamental rethinking of how knowledge is collected, represented, and accessed in OG-RAG systems \cite{carroll2020care}. Current ontology development methodologies lack frameworks for ensuring that cultural knowledge is represented authentically and ethically.

\subsubsection{Evaluation and Cultural Fidelity Assessment}

The evaluation of OG-RAG systems for culturally sensitive applications remains problematic. Traditional metrics like BLEU and ROUGE are inadequate for assessing cultural fidelity and contextual appropriateness \cite{brown2023figurative}. The development of culturally-aware evaluation frameworks requires collaboration with cultural experts and community stakeholders, but current research lacks systematic approaches for such collaboration.

Human evaluation remains the gold standard for assessing cultural fidelity, but it is expensive, time-consuming, and difficult to scale. The exploration of LLM-as-a-judge approaches shows promise but raises questions about the cultural competence of evaluation models and their ability to assess nuanced cultural appropriateness \cite{wang2024pandalm}.

\subsection{Emerging Opportunities and Breakthrough Potential}

Despite these challenges, several emerging research directions offer significant potential for advancing OG-RAG systems and addressing current limitations.

\subsubsection{Automated Ontology Learning and Community-Driven Construction}

The integration of large language models with crowdsourcing and community engagement platforms presents opportunities for scalable, participatory ontology construction. Recent work on collaborative knowledge base construction \cite{zhou2024collaborative} demonstrates that combining automated tools with community input can produce high-quality knowledge representations while reducing the burden on individual experts.

The development of user-friendly ontology editing tools powered by natural language interfaces could democratize ontology construction, enabling domain experts without formal ontological training to contribute to knowledge base development. This approach is particularly relevant for cultural heritage applications where community stakeholders should be central to the knowledge representation process.

\subsubsection{Multimodal Cultural Preservation Applications}

The integration of multimodal capabilities in OG-RAG systems opens new possibilities for comprehensive cultural preservation. The GraphVis framework \cite{zhang2024graphvis} demonstrates the potential for integrating visual and textual information, but future systems could incorporate audio, video, and other multimedia content to create rich, multifaceted representations of cultural knowledge.

The application of OG-RAG systems to oral traditions, ceremonial practices, and artisanal knowledge could provide unprecedented capabilities for preserving and transmitting cultural heritage. However, such applications require careful attention to cultural protocols and community consent processes.

\subsubsection{Cross-Lingual Ontology Transfer and Low-Resource Language Support}

The development of cross-lingual ontology transfer methods could significantly reduce the resource requirements for supporting low-resource languages. Current research on multilingual knowledge graphs \cite{chen2024multilingual} provides a foundation, but effective transfer requires preserving cultural and linguistic nuances that may not have direct equivalents across languages.

Zero-shot and few-shot learning approaches for ontology adaptation show particular promise for low-resource contexts. The ability to bootstrap comprehensive ontologies from minimal examples could make OG-RAG systems accessible for languages and cultures that currently lack sufficient digital resources.

\subsubsection{Neurobiologically-Inspired Architectures}

The HippoRAG framework \cite{bai2024hipporag} demonstrates the potential for neurobiologically-inspired approaches to knowledge organization and retrieval. Future research could explore other cognitive mechanisms, such as episodic memory consolidation and associative reasoning, to develop more sophisticated and human-like knowledge integration systems.

The integration of attention mechanisms inspired by cognitive science could enable OG-RAG systems to better model the contextual and temporal aspects of knowledge use. Such systems could adapt their retrieval and generation strategies based on the specific cognitive demands of different tasks and contexts.

\subsection{Methodological Innovations and Technical Breakthroughs}

Several technical innovations could fundamentally advance the capabilities of OG-RAG systems and address current limitations.

\subsubsection{Hierarchical and Federated Ontology Architectures}

The development of hierarchical ontology architectures that can integrate multiple levels of abstraction could address the scalability challenges of current approaches. Such architectures could support domain-specific ontologies while maintaining connections to broader foundational knowledge, enabling both specialized reasoning and cross-domain integration.

Federated ontology systems that allow distributed knowledge maintenance while preserving global coherence could address the governance and scalability challenges of large-scale ontology development. Such systems could enable community-driven knowledge development while maintaining the formal rigor required for effective reasoning.

\subsubsection{Dynamic Ontology Evolution and Learning}

Current OG-RAG systems assume static ontologies that must be manually updated. The development of systems that can automatically evolve their ontological structures based on new evidence and community input could significantly improve their adaptability and longevity. Such systems would need to balance stability with flexibility, ensuring that fundamental knowledge relationships are preserved while allowing for growth and refinement.

Machine learning approaches for ontology evolution could leverage user interactions and system performance to identify areas where ontological improvements are needed. Such adaptive systems could continuously improve their representation quality while maintaining backward compatibility with existing knowledge.

\subsection{Integration with Emerging Technologies}

The future of OG-RAG systems will likely involve integration with several emerging technology areas that could significantly enhance their capabilities.

\subsubsection{Quantum Computing for Graph Reasoning}

Quantum computing approaches to graph reasoning could provide exponential speedups for certain types of ontological reasoning tasks. While current quantum hardware is limited, the development of quantum algorithms for graph traversal and pattern matching could eventually enable OG-RAG systems to handle vastly larger and more complex knowledge graphs.

\subsubsection{Blockchain for Knowledge Provenance and Rights Management}

The integration of blockchain technologies could address some of the ethical and governance challenges in cultural knowledge representation. Blockchain-based provenance tracking could ensure that knowledge contributions are properly attributed and that cultural protocols are maintained throughout the knowledge lifecycle.

Smart contracts could automate compliance with cultural protocols and ensure that knowledge access follows community-specified guidelines. Such systems could provide the technical infrastructure needed to implement Indigenous Data Sovereignty principles in practical OG-RAG deployments.

\section{Synthesis and Implications for Future Research}

The analysis of the current state, historical evolution, and future potential of OG-RAG systems reveals several key insights that should guide future research and development efforts.

\subsection{The Paradigm Shift Toward Structured Knowledge}

The evidence presented demonstrates that the field has undergone a fundamental paradigm shift from vector-similarity-based retrieval toward structured knowledge integration. This shift is not merely technical but represents a deeper understanding that effective knowledge-grounded generation requires explicit representation of relationships, hierarchies, and domain-specific constraints.

The quantitative improvements achieved by ontology-grounded approaches—55\% improvement in factual recall and 40\% enhancement in response correctness—provide compelling evidence that this paradigm shift is not only theoretically sound but practically beneficial. These results suggest that future research should prioritize the development of increasingly sophisticated knowledge structures rather than focusing solely on scaling existing vector-based approaches.

\subsection{The Critical Role of Domain Expertise and Community Engagement}

The historical analysis reveals that successful OG-RAG systems require deep integration of domain expertise throughout the development process. The challenges faced by early knowledge integration attempts underscore the importance of involving domain experts not merely as consultants but as central participants in the knowledge representation process.

For cultural heritage applications, this insight is particularly critical. The ethical requirements for community consent and participation are not merely procedural but fundamental to the success of the technical approach. Future research must develop methodologies that genuinely center community knowledge and decision-making in the development process.

\subsection{The Need for Interdisciplinary Integration}

The convergence of multiple research streams that led to contemporary OG-RAG systems demonstrates the importance of interdisciplinary collaboration. The integration of insights from cognitive science (HippoRAG), formal ontology (OG-RAG), and graph theory (GNN-RAG) has produced systems that exceed the capabilities of any single approach.

Future breakthroughs will likely require even deeper interdisciplinary integration, incorporating insights from anthropology, linguistics, ethics, and other fields that have traditionally been separate from AI research. The development of methodologies for such integration represents a critical research challenge.

\subsection{Implications for Cultural Heritage and Low-Resource Language Applications}

The specific application to cultural heritage and low-resource languages reveals both the potential and the challenges of OG-RAG systems for preserving and transmitting cultural knowledge. The technical capabilities exist to create sophisticated knowledge representations, but the ethical and practical challenges of doing so appropriately remain significant.

The success of OG-RAG systems in addressing data scarcity through structured knowledge representation suggests that these approaches could be particularly valuable for low-resource contexts. However, the resource requirements for ontology construction may paradoxically exclude the very communities that could benefit most from these technologies.

Future research must address this tension by developing scalable, community-appropriate methodologies for ontology construction and maintenance. The integration of automated tools with community engagement processes represents a promising direction, but substantial research is needed to develop effective approaches.

\section{Conclusion}

This State-of-the-Art analysis reveals that Ontology-Grounded RAG systems represent a mature and rapidly evolving field with significant potential for advancing knowledge-intensive AI applications. The historical trajectory from traditional vector-based retrieval to sophisticated ontology-grounded approaches demonstrates the field's responsiveness to fundamental challenges and its capacity for paradigm-shifting innovation.

The contemporary landscape of OG-RAG systems offers robust technical foundations for applications requiring high factual accuracy and domain-specific reasoning. The quantitative performance improvements achieved by leading systems provide compelling evidence for the effectiveness of structured knowledge integration approaches.

However, the analysis also reveals significant challenges that must be addressed for OG-RAG systems to achieve their full potential. The resource requirements for ontology construction, the ethical challenges of cultural knowledge representation, and the limitations of current evaluation methodologies represent critical barriers to widespread adoption.

The future research directions identified in this analysis suggest that the field is poised for continued advancement, with emerging opportunities in automated ontology construction, multimodal integration, and community-driven knowledge development. The integration of OG-RAG systems with emerging technologies like quantum computing and blockchain could further expand their capabilities and address current limitations.

For the specific application to cultural heritage and low-resource languages, OG-RAG systems offer unprecedented opportunities for preserving and transmitting cultural knowledge. However, realizing this potential requires genuine engagement with ethical frameworks and community-centered development methodologies.

The three-part argument presented in this review—where we are now, how we got here, and where we could go next—demonstrates that OG-RAG systems have evolved from promising research prototypes to mature technical frameworks capable of addressing real-world challenges. The field's continued evolution will depend on its ability to integrate technical innovation with ethical responsibility and community engagement, ensuring that these powerful technologies serve the communities whose knowledge they aim to preserve and transmit.
